\documentclass[aps,prd,onecolumn,amsmath,amssymb]{revtex4-2}

\usepackage{amsmath}
\usepackage{amssymb}
\usepackage{amsthm}
\usepackage{graphicx}
\usepackage{hyperref}
\usepackage{xcolor}
\usepackage{booktabs}
\usepackage{aas_macros}

\begin{document}

\title{Spectroscopy as a Probe of Spacetime Foam}

\author{A.\ G.\ Kim}
\email{agkim@lbl.gov}
\affiliation{Physics Division, Lawrence Berkeley National Laboratory, Berkeley, California 94720, USA}
\author{P.\ E.\ Nugent}
\affiliation{Applied Mathematics and Computational Research Division, Lawrence Berkeley National Laboratory, Berkeley, California 94720, USA}
\author{L.\ Wang}
\affiliation{Department of Physics and Astronomy, 4242 TAMU, College Station, Texas 77843-4242, USA}

\date{\today}

\begin{abstract}
Quantum-gravitational spacetime foam imprints an achromatic, stochastic
arrival-time jitter on light from distant sources. We show that this jitter
modifies the first-order temporal coherence $g^{(1)}(\tau)$ and hence the
observed spectrum, providing a spectroscopic probe complementary to existing
imaging-channel (PSF, Strehl) constraints. In the
dynamical-foam limit the spectrum splits into a coherent line core suppressed
by $e^{-\sigma_\phi^2}$ and an incoherent pedestal, yielding a spectral
shape-distortion bound on the Ng--Perlman foam exponent $\alpha$.
Applied to the MeV emission lines of GRB\,221009A in the dynamical-foam
limit ($\tau_{\rm foam}\ll\tau_c$) we obtain $\alpha \gtrsim 0.63$.
For the broadband CMB the foam induces an $e^{-\sigma_\phi^2(\nu)}$ deficit
($\sigma_\phi^2 \propto \nu^2$) with a low-frequency pedestal; a Fisher
analysis of COBE/FIRAS reaches $\alpha \lesssim 0.52$ under the
physical Planck-scale pedestal kernel (essentially zero in-band;
a $\Delta\alpha\approx 0.01$ shift under an adversarial band-filling
template, modulo a $\Delta\alpha\sim 0.03$ cosmological-redshift
systematic); the cosmological-accumulation variant tightens the
FIRAS threshold to $\alpha\lesssim 0.55$, with PIXIE projected to
reach $\alpha\lesssim 0.54$ from instrumental noise improvement alone.  The final FIRAS threshold awaits
a first-principles derivation of the foam temporal autocorrelation
$C_\alpha(\tau)$, which determines the pedestal shape and hence the
foam--$\mu$ degeneracy that currently dominates the error budget. None of these bounds exclude
the holographic value $\alpha = 2/3$, which would require a confirmed
narrow spectral line at $E_{\rm obs} \gtrsim 720$\,MeV in a $z \sim 0.15$ GRB.
\end{abstract}

\maketitle

\section{Introduction}
\label{sec:intro}

A long-standing expectation from the union of quantum mechanics and general
relativity is that the smooth pseudo-Riemannian manifold of classical spacetime
breaks down at the Planck length
$\ell_{\rm P} = \sqrt{\hbar G/c^3} \simeq 1.616\times 10^{-35}$~m.
In Wheeler's original ``spacetime foam'' picture~\cite{Wheeler1955}, quantum
fluctuations of the metric on this scale render the geometry stochastic;
similar conclusions are reached in the Euclidean path-integral
treatment~\cite{Hawking1978} and in modern holographic and entropic-gravity
arguments from very different starting points.
While a complete theory of quantum gravity remains elusive, the question of
whether residual imprints of Planck-scale geometric fluctuations can accumulate
over cosmological baselines and become observable is one of the few empirical
handles on physics in the Planck regime, and it has motivated a substantial
body of phenomenological work over the past two
decades~\cite{Carlip2023,NgPerlman2022}.

A widely studied parametrization due to Ng and
collaborators~\cite{Ng2001,Ng2003,NgPerlman2022}, the $\alpha$-foam model,
encodes the foam dynamics in a power-law path-length variance
\begin{equation}
  \sigma_\ell^2(D) = A_\alpha\,D^{2(1-\alpha)}\,\ell_{\rm P}^{2\alpha},
  \label{eq:NgPerlmanVar}
\end{equation}
where $D$ is the proper distance to the source, $\alpha$ parametrizes the foam
model, and $A_\alpha$ is a dimensionless coefficient of order unity.
A photon traversing distance $D$ passes through approximately
$N \sim D/\ell_{\rm P}$ independent Planck-scale spacetime cells, each imparting
a tiny random displacement to the photon path; the exponent $\alpha$ controls
how the per-cell displacements accumulate: for random-walk foam ($\alpha=1/2$),
contributions add in quadrature ($\sigma_\ell \propto \sqrt{D\ell_{\rm P}}$);
holographic foam ($\alpha=2/3$) grows more slowly, constrained by an area-law
bound on the number of independent degrees of freedom per Planck area; and
conservative foam ($\alpha=1$) yields a distance-independent
$\sigma_\ell \propto \ell_{\rm P}$, corresponding to perfectly correlated cell
contributions.

The primary observational probe of stochastic foam to date has been the imaging
channel. Spacetime foam accumulates a cumulative phase dispersion
$\Delta\phi \simeq 2\pi\ell^{1-\alpha}\ell_{\rm P}^{\alpha}/\lambda$ across the
entrance aperture, degrading PSF sharpness and erasing interference fringes once
$\Delta\phi \gtrsim 2\pi$~\cite{LieuHillman2003}.
For PKS\,1413$+$135 ($\ell = 1.216$~Gpc) observed with HST at $\lambda=1.6\,\mu$m,
the random-walk model ($\alpha=1/2$) predicts $\Delta\phi \sim 10\times 2\pi$,
so the detection of Airy rings in this image marginally rules out that model;
the holographic model ($\alpha=2/3$) predicts $\Delta\phi \sim 10^{-9}\times 2\pi$
and is obviously not ruled out~\cite{NgPerlman2022,Christiansen2011,Perlman2015}.
Models with long-range transverse correlations evade even the random-walk
bound~\cite{LeeZurekChen2024}.

The most luminous GRB ever recorded, GRB\,221009A at $z=0.151$, has been
examined in the imaging channel by Steinbring~\cite{Steinbring2023} and
Steinbring \& Ng~\cite{SteinbringNg2026}.
Using a multiwavelength average of foam-induced angular broadening ---
analogous to a seeing disk from atmospheric turbulence ---
Steinbring~\cite{Steinbring2023} showed that the holographic ($\alpha=2/3$)
model predicts an angular widening below $\sim\!1^{\circ}$, consistent
with the Fermi and Swift localization data for GRBs at $z\leq 5$ up to
peak energies of $\sim\!10\,\mathrm{MeV}$.
Steinbring \& Ng~\cite{SteinbringNg2026} extended this to the full
multiwavelength follow-up of GRB\,221009A --- from optical and near-infrared
through X-rays and the highest-energy gamma-rays ever attributed to a GRB ---
computing the holographic-foam-widened PSF for real telescopes and finding all
data consistent with $\alpha=2/3$.
These imaging-channel analyses show that the holographic model is not
excluded by the spatial localization of GRB\,221009A.

We develop a framework based on the first-order temporal coherence
function $g^{(1)}(\tau)$ of light from astrophysical sources.
This function is complementary to the first-order spatial coherence
function $g^{(1)}(B)$, which forms the basis of imaging-channel probes.
The first-order coherence --- the field--field autocorrelation ---
inherits a foam suppression factor
that decreases continuously from unity toward complete fringe erasure as
the accumulated phase variance $\sigma_\phi^2$ grows.
In our framework the phase variance arises from temporal fluctuations of
foam, whereas in the imaging channel it arises from spatial fluctuations;
correspondingly, foam imprints itself on the spectrum rather than on the
spatial brightness profile.

The remainder of the paper is organized as follows.
Section~\ref{sec:model} introduces the spacetime foam model: the
arrival-time jitter parameterization, the temporal autocorrelation
$C_\alpha(\tau)$, and the distribution of differential jitter.
Section~\ref{sec:g1} derives the first-order coherence function
$g^{(1)}_{\rm obs}(\tau)$ of the observed field and establishes the
Gaussian closed form.
Section~\ref{sec:spec} translates these results into spectroscopic observables
via the Wiener--Khinchin theorem, establishing the split of the observed
spectrum into a coherent line core suppressed by $e^{-\sigma_\phi^2}$ and a
broad pedestal, and connects them to the imaging Strehl ratio.
Section~\ref{sec:sources} applies the spectral shape-distortion bound to two sources:
GRB\,221009A (Section~\ref{sec:grb221009a}) and the CMB as observed by
COBE/FIRAS (Section~\ref{sec:CMB}).
Section~\ref{sec:conclusions} summarizes the main results and outlines future
observational directions.

% ==============================================================
\section{Spacetime Foam Model}
\label{sec:model}
% ==============================================================

\subsection{Arrival-Time Jitter as a Time Reparameterization}

We model the detected field as a stochastically time-warped version of the
intrinsic source field:
\begin{equation}
  E_{\rm obs}(t) = E_{\rm src}\!\left(t - \delta t(t)\right),
  \label{eq:Eobs}
\end{equation}
where $E_{\rm src}(t)$ is a single spatial mode of a stationary zero-mean
complex Gaussian process (chaotic light, appropriate for the thermal and incoherent
non-thermal sources considered here, including GRB high-energy
emission and the CMB) and $\delta t(t)$ is the foam-induced
arrival-time jitter, a stationary stochastic process statistically independent
of $E_{\rm src}$.
The jitter arises from cumulative path-length fluctuations along the
line of sight of propagation distance $D$:
\begin{equation}
  \delta t(t) = \frac{1}{c}\,\Delta\ell_{\rm tot}(t).
  \label{eq:jitter}
\end{equation}
The vacuum $c$ in Eq.~\eqref{eq:jitter} reflects the standard
Ng--Perlman treatment of foam as a small perturbation atop a flat
(or FLRW) background, with photons travelling at $c$ along
stochastically perturbed geodesics.
In the geometric-optics limit the path-length shift $\Delta\ell_{\rm tot}$
is a purely geometric quantity and hence frequency-independent (achromatic),
in contrast to plasma dispersion which scales as $\lambda^2$.
This achromaticity distinguishes foam from other astrophysical screens
--- such as plasma dispersion, interstellar scintillation, and atmospheric
turbulence --- all of which imprint frequency-dependent signatures.
For broadband sources such as the CMB, the foam distortion is likewise
distinct from standard spectral distortions ($\mu$, $y$) that have
different frequency dependences. 

\subsection{Temporal Decorrelation of Spacetime Foam}
\label{sec:Calpha_origin}

The scalar variance $\sigma_\ell^2(D)$ (Eq.~\eqref{eq:NgPerlmanVar})
characterizes the amplitude of foam-induced path-length fluctuations along
a single ray; here we describe how those fluctuations are correlated
across time (temporal decorrelation).

The temporal evolution of spacetime foam --- an aspect not specified by
Eq.~\eqref{eq:NgPerlmanVar} --- determines how the arrival-time jitter
$\delta t(t)$ along a single line of sight is correlated across the lag $\tau$.
Quantum foam models posit that each Planck-scale cell along the line of sight
fluctuates on a characteristic timescale $\tau_{\rm foam}$, so that the jitter
at times $t$ and $t+\tau$ is correlated for $\tau \ll \tau_{\rm foam}$ and
essentially independent for $\tau \gg \tau_{\rm foam}$.
We introduce $C_\alpha(\tau;D)$ as the normalized \emph{temporal}
autocorrelation of the integrated path-length fluctuation
$\Delta\ell_{\rm tot}(t)$, i.e.\ $C_\alpha(\tau;D) \equiv
\langle\Delta\ell_{\rm tot}(t)\,\Delta\ell_{\rm tot}(t+\tau)\rangle
/\sigma_\ell^2(D)$, which inherits the temporal structure of the
per-cell foam dynamics over the $\sim D/\ell_{\rm P}$ cells crossed by
the line of sight.
Because $C_\alpha$ is defined on the \emph{integrated} path, its
temporal shape inherits a residual dependence on $D$ from the way
per-cell foam dynamics combine along the line of sight; the explicit
$D$-argument records this dependence.
Equivalently,
\begin{equation}
  \big\langle \Delta\ell_{\rm tot}(t)\,
              \Delta\ell_{\rm tot}(t+\tau) \big\rangle
  = \sigma_\ell^2(D)\,C_\alpha(\tau;D).
  \label{eq:pathcov1ray}
\end{equation}

The timescale $\tau_{\rm foam}$ is an additional free parameter of the
foam model, not fixed by the $\alpha$-model variance relation alone.
In the standard phenomenological framework~\cite{AmelinoCamelia2013,Ng2003},
each of the $\sim D/\ell_P$ foam cells along the line of sight
independently adopts a new configuration on timescale
$\tau_P \equiv \ell_P/c \approx 5\times10^{-44}$~s; the integrated
path-length jitter $\delta\ell_{\rm tot}(t)$ --- a sum of
$\sim D/\ell_P$ independent contributions --- likewise decorrelates on
timescale $\tau_{\rm foam}\sim\tau_P$.  With $C_\alpha(\tau)\approx 0$
for all nonzero accessible lags, the foam is fully decorrelated on any
observable timescale --- the most optimistic scenario for detection ---
and for any astrophysical source considered here
($\tau_c/\tau_{\rm foam}\gtrsim 10^{22}$ even at $37$\,MeV) the
analysis sits deep within the dynamical regime.  Models with long-range
\emph{temporal} correlations --- the temporal analogue of the
long-range \emph{transverse} spatial correlations of
Ref.~\cite{LeeZurekChen2024} that evade interferometric imaging bounds
--- would give $C_\alpha(\tau)\approx 1$ over macroscopic timescales
and thereby drive $\sigma_\Delta^2(\tau)\to 0$, suppressing
\emph{both} the GRB spectroscopic signal of
Section~\ref{sec:grb221009a} and the CMB spectral distortion of
Section~\ref{sec:CMB}.  The two channels therefore probe the same temporal
decorrelation regime, and both depend on the dynamical-foam assumption
$\tau_{\rm foam}\ll\tau_c$ adopted here.

\subsection{Differential Jitter and Its Distribution}
\label{sec:differential_jitter}

With the temporal covariance model in place, we define the key derived quantity
that enters the coherence integrals.
Define the differential jitter between photons arriving at
times $t$ and $t+\tau$:
\begin{equation}
  \Delta_t(\tau)
  \equiv \delta t(t+\tau) - \delta t(t).
  \label{eq:DeltaExt}
\end{equation}
For fixed $D$, stationarity implies that the distribution of
$\Delta_t(\tau)$ depends only on $\tau$; we suppress the $t$
subscript and write $\Delta(\tau)$ henceforth.
Expanding the square and applying Eq.~\eqref{eq:pathcov1ray}:
\begin{align}
  \sigma_\Delta^2(\tau)
  &= \frac{1}{c^2}\big\langle
    [\Delta\ell_{\rm tot}(t+\tau)
     - \Delta\ell_{\rm tot}(t)]^2
    \big\rangle \notag \\
  &= \frac{1}{c^2}\left[
    \sigma_\ell^2 + \sigma_\ell^2
    - 2\sigma_\ell^2\,C_\alpha(\tau)
    \right] \notag \\[4pt]
  &= \frac{2\sigma_\ell^2(D)}{c^2}
    \left[1 - C_\alpha(\tau;D)\right].
  \label{eq:diffVarExt}
\end{align}

By the central limit theorem (CLT), $\Delta(\tau)$
is a sum of contributions from many foam cells and is therefore Gaussian with zero mean:
\begin{equation}
  \Delta(\tau)
  \sim \mathcal{N}\!\left(0,\,\sigma_\Delta^2(\tau)\right),
  \label{eq:kernel}
\end{equation}
where $\sigma_\Delta^2$ is taken from Eq.~\eqref{eq:diffVarExt}; for
any source at $D \gtrsim 1$~kpc, the cell count
$N_{\rm cells} \sim D/\ell_{\rm P} \gtrsim 10^{54}$ amply satisfies
the large-$N$ requirement and higher-cumulant corrections are negligible.  For $\alpha = 1/2$ (random walk), per-cell path-length variances
are equal ($\ell_{\rm P}^2$ each) and the CLT applies in its standard form.
For $\alpha \neq 1/2$, the per-cell contributions scale differently with $D$
and the equality of per-cell variances is a model assumption rather than a
direct CLT consequence; Gaussianity is adopted as the minimal assumption
consistent with the $\alpha$-model.

The characteristic function of $\Delta(\tau)$ at carrier frequency
$\omega_0$ defines a foam suppression factor applicable to any source
profile.  Since $\Delta(\tau)$ is Gaussian,
\begin{equation}
  \big\langle e^{i\omega_0\Delta(\tau)}\big\rangle_{\delta t}
  = e^{-\omega_0^2\sigma_\Delta^2(\tau)/2}
  \;\equiv\; \mathcal{F}_{\rm foam}(\tau),
  \label{eq:gaussian_jitter}
\end{equation}
with $\mathcal{F}_{\rm foam}(0) = 1$.

From the foam path-length variance $\sigma_\ell^2(D_C)$ and carrier
frequency $\omega_0$, we define $\sigma_\phi^2 \equiv
\tfrac{1}{2}\omega_0^2\cdot(2\sigma_\ell^2/c^2) = (\omega_0\sigma_\ell/c)^2$
as the asymptotic phase variance --- the value
$\tfrac{1}{2}\omega_0^2\sigma_\Delta^2(\tau)$ approaches as
$C_\alpha(\tau;D)\to 0$.  In the dynamical-foam limit, the
Ng--Perlman path-length model~\eqref{eq:NgPerlmanVar} with $D_C$ the
line-of-sight comoving distance~\cite{NgPerlman2022} gives
\begin{equation}
  \sigma_\phi^2 = \left(\frac{2\pi}{\lambda_{\rm obs}}\right)^{\!2}
    A_\alpha\,D_C^{2(1-\alpha)}\,\ell_P^{2\alpha}.
  \label{eq:sigmaphi}
\end{equation}
Substituting $\sigma_\Delta^2(\tau) =
(2\sigma_\ell^2/c^2)[1-C_\alpha(\tau;D)]$ from
Eq.~\eqref{eq:diffVarExt}, Eq.~\eqref{eq:gaussian_jitter} takes the compact form
\begin{equation}
  \mathcal{F}_{\rm foam}(\tau) = e^{-\sigma_\phi^2[1-C_\alpha(\tau;D)]}.
  \label{eq:Ffoam}
\end{equation}
At $\tau\to\infty$: $\mathcal{F}_{\rm foam}(\infty)$ equals
$e^{-\sigma_\phi^2}$ in the dynamical limit ($C_\alpha\to 0$) and
approaches unity for quasi-static foam ($C_\alpha\approx 1$).
Equation~\eqref{eq:sigmaphi} uses observer-frame wavelength
$\lambda_{\rm obs}$ throughout and so neglects the cosmological
redshift of the photon along the propagation path: at redshift $z'$
the photon's local wavelength was $\lambda(z') = \lambda_{\rm obs}/(1+z')$,
and the phase accumulated per unit path length scales as $1/\lambda^2$.
For a model in which foam path-length variance accumulates uniformly in
comoving distance, the redshift-weighted form is
\begin{equation}
  \sigma_\phi^2
  \;=\; \left(\frac{2\pi}{\lambda_{\rm obs}}\right)^{\!2}
        \int_0^{z_{\rm s}}\!(1+z')^2\,\frac{d\sigma_\ell^2}{dD_C}\,
        \frac{dD_C}{dz'}\,dz'\,,
  \label{eq:sigmaphi_cosmo}
\end{equation}
which reduces to Eq.~\eqref{eq:sigmaphi} when $(1+z')^2$ can be taken
constant over the path.  For the GRB at $z=0.151$ the path-averaged
$\langle(1+z')^2\rangle_{D_C}\approx 1.16$, enhancing $\sigma_\phi^2$
by $16\%$ --- negligible at the precision of the present analysis.
For the CMB at $z_{\rm LS}\approx 1089$ the weighting is dominated by
the high-$z$ portion of the path: with Planck~2018 flat $\Lambda$CDM
the size and sign of this correction depend on how foam fluctuations
evolve in cosmic time; two natural anchorings bracket the range:
(a)~variance accumulating uniformly in comoving distance,
$d\sigma_\ell^2/dD_C = \text{const}$, enhances $\sigma_\phi^2$ by a
factor of $\sim 10^4$;
(b)~variance anchored to proper-now coordinates with foam frozen at
the present epoch gives $(1+z')^2\to 1$ inside the integral,
leaving $\sigma_\phi^2$ unchanged.
Bounds quoted below use Eq.~\eqref{eq:sigmaphi} with $A_\alpha=1$ and
adopt anchoring (b); the CMB systematic of Section~\ref{sec:CMB}
discusses the corresponding shift under anchoring (a).

% ==============================================================
\section{First-Order Coherence }
\label{sec:g1}
% ==============================================================

\subsection{Jitter Convolution Formula for $g^{(1)}$}

The observed first-order coherence is
\begin{equation}
  g^{(1)}_{\rm obs}(\tau) = \frac{\langle E^*_{\rm obs}(t)\,E_{\rm obs}(t+\tau)\rangle}
                                  {\langle|E_{\rm obs}(t)|^2\rangle}.
  \label{eq:g1def}
\end{equation}
Substituting Eq.~\eqref{eq:Eobs} and letting
$u = t - \delta t(t)$, $v = t+\tau - \delta t(t+\tau)$, so that
$v - u = \tau - \Delta(\tau)$, stationarity of $E_{\rm src}$ gives
\begin{equation}
  \langle E^*_{\rm src}(u)\,E_{\rm src}(v)\rangle
  = \langle|E_{\rm src}|^2\rangle\,g^{(1)}_{\rm src}(\tau - \Delta(\tau)).
\end{equation}
Since $\delta t$ is independent of $E_{\rm src}$, and since
$\langle|E_{\rm obs}(t)|^2\rangle = \langle|E_{\rm src}(t-\delta t)|^2\rangle
= \langle|E_{\rm src}|^2\rangle$ by stationarity (so the normalisation
in Eq.~\eqref{eq:g1def} cancels), averaging over the jitter yields
\begin{equation}
  g^{(1)}_{\rm obs}(\tau)
    = \big\langle g^{(1)}_{\rm src}(\tau - \Delta(\tau))\big\rangle_{\delta t}
    = \int_{-\infty}^{\infty} \mathcal{N}(\Delta;0,\sigma_\Delta^2(\tau))\,
      g^{(1)}_{\rm src}(\tau - \Delta)\,d\Delta.
  \label{eq:g1obs}
\end{equation}
The observed first-order coherence is the intrinsic coherence function
convolved with the Gaussian differential-jitter kernel.

\subsection{Gaussian Source Profile}
\label{sec:closedforms}

The GRB\,221009A emission line is well described by a Gaussian profile
in each epoch~\cite{YQZhang2024}, and it serves as the representative
case admitting the simplest closed form.
A Gaussian line of $1\sigma$ width $\sigma_\nu$ centred on $\nu_0$
corresponds via the Wiener--Khinchin theorem to the coherence function
\begin{equation}
  g^{(1)}_{\rm src}(\tau) \;=\; e^{-i\omega_0\tau}\,e^{-\tau^2/2\tau_G^2},
  \label{eq:g1src_gauss}
\end{equation}
where $\omega_0 \equiv 2\pi\nu_0$ and $\tau_G \equiv 1/(2\pi\sigma_\nu)$ is the
Gaussian coherence decay time.  The carrier factor $e^{-i\omega_0\tau}$
encodes the line's position in the spectrum and must be retained: dropping
it would substitute a baseband Gaussian centred at $\nu=0$ for the physical
line and would lose the dominant foam effect (line-peak suppression).
Substituting into the jitter-convolution integral~\eqref{eq:g1obs} and
completing the square gives the exact closed form
\begin{align}
  g^{(1)}_{\rm obs}(\tau)
  &= \frac{\tau_G}{\sqrt{\tau_G^2+\sigma_\Delta^2(\tau)}}\,
     \exp\!\left[\,-\,i\,\omega_0\tau\,
       \frac{\tau_G^2}{\tau_G^2+\sigma_\Delta^2(\tau)}\right] \notag\\
  &\quad\times \exp\!\left[\,-\,\frac{\tau^2+\omega_0^2\,
       \tau_G^2\,\sigma_\Delta^2(\tau)}
       {2\bigl[\tau_G^2+\sigma_\Delta^2(\tau)\bigr]}\right].
  \label{eq:g1obs_gauss}
\end{align}
The observed coherence is again a Gaussian-modulated oscillation, but
with three foam-induced modifications: (i) a slight redshift of the
carrier, $\nu_0^{\rm eff} = \nu_0\,\tau_G^2/[\tau_G^2+\sigma_\Delta^2(\tau)]$;
(ii) a wider Gaussian decay of $g^{(1)}_{\rm obs}(\tau)$ in lag (decay
scale $\sqrt{\tau_G^2+\sigma_\Delta^2(\tau)}$ vs.\ $\tau_G$);
and (iii) a global amplitude suppression by the factor
$\exp\!\left[-\omega_0^2\tau_G^2\sigma_\Delta^2/\{2(\tau_G^2+\sigma_\Delta^2)\}\right]$.
Normalisation $g^{(1)}_{\rm obs}(0)=1$ holds because $\sigma_\Delta(0)=0$.
The narrowband limit $\tau_G\to\infty$ of
Eq.~\eqref{eq:g1obs_gauss} reduces to a pure phase-modulated
oscillation, which is the single-mode building block from which the
broadband CMB spectrum of Section~\ref{sec:CMB_derivation} is
constructed by incoherent superposition.

% ==============================================================
\section{Implications for Spectroscopy}
\label{sec:spec}
% ==============================================================

\subsection{The Wiener--Khinchin Bridge}

The foam-modified observed coherence function — given in general by the
jitter-convolution integral Eq.~\eqref{eq:g1obs} with closed forms
Eq.~\eqref{eq:g1obs_gauss} — carries direct
spectroscopic content: for a stationary field the power spectral density
$S(\nu)$ and $g^{(1)}(\tau)$ are Fourier-transform pairs,
implying
\begin{equation}
  S_{\rm obs}(\nu) = \mathcal{F}\!\left\{g^{(1)}_{\rm obs}(\tau)\right\}(\nu).
  \label{eq:WK_obs}
\end{equation}
The theorem applies because astronomical emission lines have finite
linewidth (so the Fourier transform is well defined) and because the
foam-induced jitter is a stationary process independent of the source
(foam acts on photons after emission, while $\delta t(t)$ statistics are
time-translation invariant on observational timescales).  The observed
field is then stationary when the integration time exceeds the foam
coherence time $\tau_{\rm foam}$ (the lag scale beyond which
$C_\alpha(\tau;D)\to 0$), so that the time-averaged spectrum coincides
with its ensemble average: $T_{\rm int}\gg\tau_{\rm foam}$.

In the dynamical limit $C_\alpha(\tau;D)\to 0$ for
$|\tau|\gtrsim\tau_{\rm foam}$, $\mathcal{F}_{\rm foam}(\tau)$ acquires
a long-lag plateau $e^{-\sigma_\phi^2}$ outside a narrow window
$|\tau|\lesssim\tau_{\rm foam}$.  Fourier-transforming via
Eq.~\eqref{eq:WK_obs}, the plateau produces a coherent core
$e^{-\sigma_\phi^2}S_{\rm src}(\nu)$, with the complementary fraction
$1-e^{-\sigma_\phi^2}$ of the source flux in a pedestal sourced by the
short-lag window.  The per-mode derivation in
Section~\ref{sec:CMB_derivation} makes this decomposition explicit for
a broadband source.

For the Gaussian source profile of Section~\ref{sec:closedforms},
Fourier-transforming Eq.~\eqref{eq:g1obs_gauss} with $\sigma_\Delta^2$
held at its asymptotic value $2\sigma_\ell^2/c^2$ gives the explicit
Gaussian core
\begin{equation}
  S_{\rm obs}^{\rm core}(\nu) \;=\;
    \frac{\tau_G}{\sqrt{\tau_G^2+\sigma_\Delta^2}}\,
    \exp\!\left[-\frac{\sigma_\phi^2\,\tau_G^2}{\tau_G^2+\sigma_\Delta^2}\right]
    \cdot
    \frac{1}{\sqrt{2\pi}\,\sigma_{\rm obs}}
    \exp\!\left[-\frac{(\nu-\nu_0^{\rm eff})^2}{2\sigma_{\rm obs}^2}\right],
  \label{eq:Sobs_gauss}
\end{equation}
with $\nu_0^{\rm eff}=\nu_0\tau_G^2/(\tau_G^2+\sigma_\Delta^2)$ and
$\sigma_{\rm obs}=1/[2\pi\sqrt{\tau_G^2+\sigma_\Delta^2}]<\sigma_\nu$.
The Gaussian's carrier shift, line narrowing, and peak suppression are
the frequency-domain images of items~(i)--(iii) of
Section~\ref{sec:closedforms}.

The pedestal in either decomposition is sourced by the short-lag
window $|\tau|\lesssim\tau_{\rm foam}$ and so spreads over a frequency
range $\sim 1/\tau_{\rm foam}$.  When $\tau_{\rm foam}\ll
\tau_G,\sigma_\Delta$ the pedestal is correspondingly broad and
shallow, and the coherent core is the operative spectrum near the line.


% ==============================================================
\subsection{Connection to the Imaging Strehl Ratio}
\label{sec:connection}

The first-order coherence function $g^{(1)}_{\rm obs}(\tau)$ derived
above describes \emph{temporal} coherence along a single line of
sight, with the foam suppression governed by the temporal
decorrelation $C_\alpha(\tau;D)$ of Eq.~\eqref{eq:Ffoam}.  The Strehl
ratio of an imaging telescope is the \emph{spatial} counterpart:
phase variations across the aperture plane are governed by an
aperture-plane coherence function $C_\alpha^{\rm sp}(B;D)$, constructed
by direct analogy to the temporal $C_\alpha(\tau;D)$ of
Section~\ref{sec:Calpha_origin} under the same foam-cell statistics.

The structural equivalence is not merely formal: in both cases the
foam suppression arises as the characteristic function of a zero-mean
Gaussian differential random variable.  Temporally, the variable is
the differential jitter $\Delta(\tau)=\delta t(t+\tau)-\delta t(t)$
with variance $\sigma_\Delta^2(\tau)=(2\sigma_\ell^2/c^2)[1-C_\alpha(\tau;D)]$
(Eq.~\eqref{eq:diffVarExt}).  Spatially, the variable is the
differential phase
$\Delta\phi(B)=\phi(\mathbf{x}+\mathbf{B})-\phi(\mathbf{x})$
between two points on the aperture plane separated by baseline $B$,
with variance
\begin{equation}
  \sigma_{\Delta\phi}^2(B) \;=\; 2\sigma_\phi^2\bigl[1-C_\alpha^{\rm sp}(B;D)\bigr].
  \label{eq:dphi-variance}
\end{equation}

The focal-plane field of a telescope is the Fraunhofer Fourier
transform of the aperture-plane field, giving the diffraction-limited
PSF when the aperture phase is uniform.  Foam imprints a Gaussian
random phase $\phi(\mathbf{x})$ on the aperture with variance
$\sigma_\phi^2$ and two-point correlation $C_\alpha^{\rm sp}(B;D)$,
blurring this PSF.  In the resolved limit
$C_\alpha^{\rm sp}(D_{\rm tel};D)\to 0$ --- where aperture phases at
distinct points decorrelate --- the on-axis peak is attenuated by the
Strehl ratio
\begin{equation}
  S \;\simeq\; e^{-\sigma_\phi^2}
  \label{eq:strehl}
\end{equation}
(extended-Mar\'echal~\cite{NgPerlman2022}), the surrounding Airy
fringes are washed out by the same factor, and the complementary
$1-e^{-\sigma_\phi^2}$ of the source flux scatters into an incoherent
halo.

Imaging and spectroscopic channels each probe $\sigma_\phi^2$ through
a Fourier transform --- aperture position $\leftrightarrow$ focal-plane
angle in imaging (Fraunhofer), lag $\leftrightarrow$ frequency in
spectroscopy (Wiener--Khinchin) --- with the rapid chromatic tightening
$\sigma_\phi\propto 1/\lambda$ exploited in both Ng \&
Perlman's~\cite{NgPerlman2022} Table~1 and the analyses of
Section~\ref{sec:grb221009a}.  The shared $\sigma_\phi^2$ that
suppresses the coherent diffraction-limited PSF in imaging also
suppresses the coherent line core in spectroscopy, placing the
published imaging
constraints~\cite{NgPerlman2022,LieuHillman2003,Christiansen2011,Perlman2015}
on a common footing with the present framework.  The duality is not
perfect at the observational interface: isolating the foam contribution
requires an independently characterized telescope aperture in the
imaging channel and an independently characterized intrinsic source
spectrum in the spectroscopic channel.

% ==============================================================
\section{Sources}
\label{sec:sources}

Time-varying spacetime foam modifies the source first-order coherence,
and hence its spectrum, in a wavelength-dependent way through
the scaling $\sigma_\phi \propto 1/\lambda$. Two observable
consequences follow. First, individual spectral features are distorted according to the
Gaussian-source result Eq.~\eqref{eq:Sobs_gauss}, with a carrier shift,
line narrowing, and peak suppression that scale chromatically as
$\sigma_\phi\propto1/\lambda$. Second, the
broadband spectral energy distribution is distorted: the cosmic
microwave background, expected to be a near-perfect blackbody, would
acquire foam-induced deviations whose chromatic signature is distinct
from standard chemical-potential and Compton-$y$ distortions.
In this section we consider these two possibilities.

\subsection{GRB\,221009A}
\label{sec:grb221009a}

For spectroscopy of a single source without an absolute flux reference,
the smooth suppression factor $e^{-\sigma_\phi^2(\nu)}$ attenuates line
and continuum equally at any given frequency, preserving the equivalent
width: uniform dimming is undetectable from the spectrum alone.  The
carrier-included closed form Eq.~\eqref{eq:g1obs_gauss} supplies the
needed shape-distortion criterion: foam broadens the observed line width
by a calculable energy-dependent amount (derived below as
Eq.~\eqref{eq:robs_decomp}), and multi-epoch line-width measurements
across a wide energy range give additional leverage, as we now show
for GRB\,221009A.

GRB\,221009A (``the BOAT,'' Brightest Of All Time) is the most
energetic gamma-ray burst in the instrumental record, with isotropic-equivalent
energy $E_{\rm iso}\sim10^{55}\,\mathrm{erg}$ at redshift $z=0.151$,
a once-in-$\sim10^4$-year event~\cite{Burns2023}.
Two independent spectral analyses have reported a narrow gamma-ray
emission feature in the Fermi/GBM data, interpreted as blueshifted
electron--positron pair annihilation~\cite{Ravasio2024,YQZhang2024}, at
observed energies spanning $6$--$37\,\mathrm{MeV}$.
This pair-annihilation identification is one proposed interpretation
and is not yet definitive; the foam analysis depends only on
Doppler-invariance of the line's fractional width, not on the specific
emission mechanism.
We exploit the \emph{multi-epoch} constancy of the fractional line
width across this factor-of-six energy range to obtain a stringent
constraint.
Adopting flat $\Lambda$CDM
($\Omega_M=0.27$, $\Omega_\Lambda=0.73$,
$H_0=70\,\mathrm{km\,s^{-1}\,Mpc^{-1}}$), the comoving distance is
$D_C(z=0.151)\approx0.63\,\mathrm{Gpc}\approx1.93\times10^{25}\,\mathrm{m}$
and $\ln(D_C/\ell_P)\approx138.3$.

The analysis below operates in the dynamical-foam regime
$\tau_{\rm foam}\ll\tau_c$ adopted throughout
Section~\ref{sec:model}; in the opposite quasi-static limit the
framework yields no foam imprint on the line spectrum and the bound
does not apply (caveat~(vi) below).

\textit{Multi-epoch line data.}
Zhang et al.~\cite{YQZhang2024} perform a joint analysis of GECAM-C and
Fermi/GBM data covering the full prompt emission, including the initial
bright phase that saturated GBM alone.  They detect a Gaussian emission
line with centroid evolving as $E_{\rm line}\propto t^{-1}$ from
$\approx37$\,MeV down to $\approx6$\,MeV over $T_{\rm ref}+246$--$360$\,s.
In the seven sub-intervals where the line is sufficiently bright to
constrain its width, the measured fractional width
$r_{\rm obs}\equiv\sigma_{\rm obs}/E_{\rm line} = 0.10\pm0.01$ is
consistent with a single constant across the full dynamic range
($\chi^2/\mathrm{dof}=4/7$; Zhang et al.\ \S4.4 and their Figure~7c).
This constancy is the predicted behaviour of a Doppler-boosted line.
In the comoving frame of the emitting plasma the line has fractional
width $r' = \sigma'/E'$, fixed by the comoving emission physics
(pair-plasma temperature, velocity dispersion).  When the outflow has
bulk Doppler factor $\mathcal{D}$ --- epoch-dependent as the jet
decelerates --- both the line centroid and its width are multiplied by
the same factor: $E = \mathcal{D}E'$, $\sigma_{\rm int} =
\mathcal{D}\sigma'$, so $r_{\rm int} \equiv \sigma_{\rm int}/E = r'$.
The Doppler factor cancels in the ratio; $r_{\rm int}$ is a
Lorentz-invariant property of the line, independent of epoch.  The
observed $E_{\rm line}\propto t^{-1}$ centroid evolution over a factor
of $\approx 6$ in energy traces the bulk-Doppler deceleration of the
jet, and the simultaneous constancy of $r_{\rm obs}$ confirms the
prediction.  The line is independently detected by Ravasio et
al.~\cite{Ravasio2024} in Fermi/GBM data alone, with
$\sigma_{\rm gauss}=1.31^{+0.31}_{-0.30}$\,MeV at
$E_{\rm line}=12.56^{+0.30}_{-0.31}$\,MeV in $T+280$--$300$\,s
($\sigma/E\approx 10.4\%$, consistent with Zhang's
$r_{\rm obs}=0.10\pm0.01$ at $0.4\sigma$ in the time interval that
overlaps Zhang's bright-line window).  Ravasio also reports the line
in three later intervals ($T+300$--$360$\,s) where $\sigma/E$ rises
to $\approx 15$--$20\%$, so the constancy of $r_{\rm obs}$ across the
full $9$--$37$\,MeV lever arm rests on the GECAM-C$+$Fermi/GBM joint
analysis~\cite{YQZhang2024}, with Ravasio providing an
instrument-independent confirmation of the line at the
highest-signal-to-noise epoch.

\textit{Width decomposition.}
Because $r_{\rm int}$ is the same in every epoch, the multi-epoch
measurements constrain a single common $r_{\rm int}$ alongside the
foam amplitude $r_{\rm foam,0}$, giving the lever arm its statistical
power.  The Gaussian-source spectrum Eq.~\eqref{eq:Sobs_gauss} gives the
observed centroid and absolute width
$\nu_0^{\rm eff}=\nu_0\,\tau_G^2/(\tau_G^2+\sigma_\Delta^2)$ and
$\sigma_{\rm obs}=1/[2\pi\sqrt{\tau_G^2+\sigma_\Delta^2}]$.  Although
$\sigma_{\rm obs}$ is slightly \emph{narrower} than the intrinsic
$\sigma_\nu$ by a factor $\tau_G/\sqrt{\tau_G^2+\sigma_\Delta^2}$,
the centroid is redshifted by the \emph{square} of that factor, so
the fractional width that Zhang et al.\ report as $\sigma/E$ in each
epoch \emph{increases}:
\begin{equation}
  r_{\rm obs}^2(E) \;=\; r_{\rm int}^2\bigl[1+\sigma_\Delta^2/\tau_G^2(E)\bigr]
                  \;=\; r_{\rm int}^2 + r_{\rm foam}^2(E),
  \label{eq:robs_decomp}
\end{equation}
with foam excess $r_{\rm foam}^2(E)\equiv r_{\rm int}^2\,
\sigma_\Delta^2/\tau_G^2(E)$.  The relation is multiplicative in
$(1+\sigma_\Delta^2/\tau_G^2)$, not additive in independent
foam/intrinsic widths: the foam acts on the source through the
jitter-convolution Eq.~\eqref{eq:g1obs}, a \emph{multiplication} of
the spectrum by a Gaussian factor rather than a convolution with it
(Section~\ref{sec:closedforms}).

In the dynamical-foam limit ($\tau_{\rm foam}\ll\tau_c$),
$\sigma_\Delta^2(\tau)$ saturates to $2\sigma_\ell^2/c^2$ on lags
$|\tau|\gtrsim\tau_{\rm foam}$.  The broadening term in
Eq.~\eqref{eq:robs_decomp} is set by this variance at the source
coherence lag $\tau\sim\tau_G$, which satisfies
$\tau_G\sim\tau_c\gg\tau_{\rm foam}$ in this regime, so the saturated
value is the operative one.
The dominant foam effect on the spectrum is then carrier-phase
suppression of the line peak by $e^{-\sigma_\phi^2(E)}$, together with
line-width broadening captured by item~(ii) of
Section~\ref{sec:closedforms}; both follow from
Eq.~\eqref{eq:g1obs_gauss}.  Substituting
$\sigma_\Delta=\sqrt{2}\,\sigma_\ell/c$ and
$\tau_G(E)=\lambda_{\rm obs}(E)/(2\pi c\,r_{\rm int})$, the
foam-to-intrinsic timescale ratio takes the compact form
\begin{equation}
  \frac{\sigma_\Delta}{\tau_G(E)}
  \;=\; \sqrt{2}\,r_{\rm int}\,\sigma_\phi(E),
  \label{eq:sigmaDelta_tauG_dyn}
\end{equation}
the product of the intrinsic fractional width $r_{\rm int}$
(a source property) and the asymptotic phase variance
$\sigma_\phi$ (Eq.~\eqref{eq:sigmaphi}).  Substituting into
Eq.~\eqref{eq:robs_decomp} gives
$r_{\rm obs}^2(E) = r_{\rm int}^2\bigl[1+2\,r_{\rm int}^2\sigma_\phi^2(E)\bigr]$,
so the foam excess is
\begin{equation}
  r_{\rm foam}^{\rm (dyn)}(E)
  \;=\; \sqrt{2}\,r_{\rm int}^2\,\sigma_\phi(E)
  \;=\; r_{\rm int}^2\cdot\frac{2\sqrt{2}\,\pi\,\sigma_\ell\,E}{hc}
  \;\propto\; E.
  \label{eq:rfoam_dyn}
\end{equation}
The $r_{\rm int}^2$ prefactor reflects two geometric factors: one
converts the absolute foam-induced shift ($\propto\sigma_\phi\,\nu_0$)
into a fractional width; the other enters through
$\tau_G\propto 1/r_{\rm int}$ in the broadening term
$\sigma_\Delta^2/\tau_G^2$.  A narrower intrinsic line (smaller
$r_{\rm int}$, longer $\tau_G$) is therefore more robust to dynamical
foam: $\sigma_\Delta$ falls further below the source coherence time.

The fit-ready two-parameter model is
\begin{equation}
  r_{\rm obs}^2(E) \;=\; r_{\rm int}^2 + r_{\rm foam,0}^2\,(E/E_0)^2,
  \label{eq:robs_model}
\end{equation}
with $r_{\rm foam,0}\equiv r_{\rm foam}(E_0)$ the foam contribution at
the reference energy $E_0 = 37\,\mathrm{MeV}$
($\lambda_0 = 3.4\times10^{-14}\,\mathrm{m}$).

\textit{Fit and upper limit on $r_{\rm foam,0}$.}
A $\chi^2$ fit of the seven Zhang et al.\ measurements
$(E_i, r_i\pm\delta r_i)$ to Eq.~\eqref{eq:robs_model} with free
parameters $(r_{\rm int},\,r_{\rm foam,0}\geq 0)$ yields a best
fit at $r_{\rm foam,0}=0$, $r_{\rm int}=0.111$, with
$\chi^2_{\rm min}=5.3$ ($\chi^2/\mathrm{dof}=0.88$ with 6 effective
degrees of freedom; since $r_{\rm foam,0}=0$ at the boundary the fit
reduces to one free parameter $r_{\rm int}$), confirming the
constant-$r_{\rm obs}$ hypothesis adopted by Zhang et al.\ (their
\S4.4 and Figure~7c).  Profiling over $r_{\rm int}$ gives the one-sided
95\%~CL upper limit $r_{\rm foam,0}^{\rm max} = 0.27$ at
$\Delta\chi^2 = 2.71$.\footnote{The standard Wilks asymptotic result
$\Delta\chi^2\sim\chi^2(1)$ for one extra parameter assumes the
maximum-likelihood estimator lies in the \emph{interior} of the
parameter space; here the physical constraint $r_{\rm foam,0}\geq 0$
places the best fit on the boundary, so this assumption fails.  The
Chernoff (1954) boundary theorem~\cite{Chernoff1954} replaces it: when
the true value of a non-negative parameter lies on the boundary, the
asymptotic distribution of the likelihood-ratio statistic
$\Delta\chi^2$ is a 50:50 mixture of a point mass at zero and a
$\chi^2(1)$ distribution.  Heuristically, half the realizations of the
noise drive the unconstrained estimator above zero (contributing the
$\chi^2(1)$ piece), and half drive it below zero where the constraint
pins $r_{\rm foam,0}=0$ exactly (the point mass).  The 95th percentile
of the mixture therefore coincides with the 90th percentile of
$\chi^2(1)$, $\Delta\chi^2=2.71$ --- identical to the standard
one-sided Gaussian $95\%$ CL threshold, but derived for the correct
boundary-constrained sampling distribution.}  The constraint is comparatively weak
because the foam signal peaks at $E_0 = 37$\,MeV, where the single
most uncertain data point (G1) resides.  With seven measurements and
per-point fractional uncertainties of $30$--$50\%$, the
$\Delta\chi^2=2.71$ coverage relies on the Gaussian-$\chi^2$
approximation; a bootstrap or profile-likelihood Monte Carlo would
refine the exact coverage but is not expected to shift the bound by
more than a few percent.

\textit{Constraint on $\alpha$.}
Eq.~\eqref{eq:rfoam_dyn} gives
$\sigma_\ell = r_{\rm foam,0}\,\lambda_0/(2\sqrt{2}\,\pi\,r_{\rm int}^2)$;
converting via $\sigma_\ell^2 = D_C^{2(1-\alpha)}\ell_P^{2\alpha}$
(Eq.~\eqref{eq:NgPerlmanVar} with $A_\alpha=1$) gives
\begin{equation}
  \alpha \;\gtrsim\;
  \frac{\ln(D_C/\lambda_0) + |\!\ln r_{\rm foam,0}^{\rm max}|
        - |\!\ln(2\sqrt{2}\,\pi\,r_{\rm int}^2)|}
       {\ln(D_C/\ell_P)}
  \;=\; \frac{89.2 + 1.3 - 3.3}{138.3}
  \;\approx\; 0.63\,,
  \label{eq:alpha_grb_dyn}
\end{equation}
where $|\!\ln(0.27)|\approx1.3$ and
$|\!\ln(2\sqrt{2}\,\pi\,r_{\rm int}^2)|\approx 3.3$ for
$r_{\rm int}=0.066$, the extreme value within the $95\%$~CL
contour.
The cosmological redshift correction of Section~\ref{sec:closedforms}
enhances $\sigma_\phi^2$ by $16\%$ at $z=0.151$, shifting the bound
by $\Delta\alpha\sim 5\times10^{-4}$ --- negligible at the current
precision.
Using $\sigma_\Delta=\sqrt{2}\,\sigma_\ell/c$ and
$c\tau_G(E_0)=\hbar c/(r_{\rm int}E_0)\approx 8.1\times10^{-14}\,\mathrm{m}$
for $r_{\rm int}=0.066$, the ratio $\sigma_\Delta/\tau_G(E_0)\approx 4.1$
at the bound ($\sigma_\ell\approx 2.3\times 10^{-13}$~m), placing
the bound deep in the strong-foam regime $\sigma_\Delta\gtrsim\tau_G$.
The individual line-width measurements carry $30$--$50\%$ fractional
uncertainties; the multi-epoch lever arm is essential for reaching
this precision.

The G1 interval ($T+246$--$256$\,s, $E\approx37$\,MeV) lies at or
beyond the upper boundary of the GBM BGO detector range
($4$--$35$\,MeV), so its line-width measurement carries larger
systematic uncertainty than the later intervals.  Refitting the
remaining six intervals gives a best fit at $r_{\rm foam,0}=0$,
$r_{\rm int}=0.109$, with $\chi^2_{\rm min}=4.4$ over five effective
degrees of freedom, and a $95\%$~CL profile boundary at
$r_{\rm foam,0}^{\rm max}=0.27$ paired with $r_{\rm int}=0.066$ ---
yielding $\alpha\gtrsim 0.63$, statistically indistinguishable from
the all-seven-interval bound.  G1's large fractional uncertainty
($\sim 40\%$) gives it little weight in the $\chi^2$ despite its
position at the top of the lever arm.

The bound corresponds to $\sigma_\Delta\gtrsim\tau_G$,
where Eq.~\eqref{eq:g1obs_gauss} predicts \emph{two} observable
effects: (a)~line-width broadening
$r_{\rm obs}^2 = r_{\rm int}^2(1+\sigma_\Delta^2/\tau_G^2)$ --- which the
fit constrains via $r_{\rm foam,0}^{\rm max}=0.27$ --- and
(b)~a line-center redshift $\nu_0^{\rm eff}/\nu_0 =
\tau_G^2/(\tau_G^2+\sigma_\Delta^2)$ that becomes substantial at the
upper limit ($\nu_0^{\rm eff}/\nu_0\approx 0.06$ at
$\sigma_\Delta/\tau_G\approx 4.1$).  The lever-arm fit uses observed line
centers as if they were intrinsic (pair-annihilation Doppler-boosted to
$E_i$ in each epoch) and does not separately constrain (b).  Closing
the gap requires either~(i) an independent constraint on the per-epoch
bulk Doppler factor $\mathcal{D}_i$ --- from afterglow blastwave
modeling, TeV-photon photosphere arguments, or variability-timescale
estimates --- that pins the pre-foam centroid, or~(ii) detection of a
second emission line of known rest-frame energy whose distinct
$\tau_G(E)$ would produce a different foam-redshift factor in the same
epoch and thereby separate $\mathcal{D}_i$ from the foam shift.
Neither is available in the present dataset.  The bound is therefore
weaker than what a joint fit would yield, and at the upper limit it is
moreover not internally self-consistent: with $\sigma_\Delta^2/\tau_G^2
\propto E^2$, the foam-induced centroid factor $\nu_0^{\rm eff}/\nu_0$
would vary across the data range from $\approx 0.06$ at
$E=37$\,MeV to $\approx 0.7$ at $E=6$\,MeV, producing a strongly
energy-dependent depression of the observed centroid that is visibly
inconsistent with the smooth $E_{\rm line}\propto t^{-1}$
Doppler-deceleration trend in the data.  A joint center-plus-width fit
would quantify this and likely tighten the bound substantially.

If the intrinsic fractional width $r_{\rm int}$ evolves with
energy (e.g.\ due to changing pair-plasma temperature across epochs),
this could partially mimic or mask a foam signal.  Marginalizing over a
free linear-in-$E$ nuisance term $r_{\rm int}(E) = r_0 + r_1(E/E_0)$
leaves the bound essentially unchanged
($\alpha\gtrsim 0.63$), because the $E^2$ foam scaling has
distinct curvature from a linear intrinsic trend across the
factor-of-six energy range.  Numerically, the best-fit linear slope
$r_1$ is consistent with zero and, although the 95\%~CL upper limit
on $r_{\rm foam,0}$ broadens from $0.27$ to $0.31$ when the nuisance
is included, the corresponding conservative $\alpha$ bound shifts by
$\Delta\alpha\approx 3\times 10^{-3}$ --- a direct consequence of the
logarithmic compression in the Ng--Perlman
relation --- confirming that the $E^2$ and linear-$E$ shapes are
effectively orthogonal for the seven-point lever-arm spanning
$9$--$37$~MeV.

% ==============================================================
\subsection{CMB}
\label{sec:CMB}

For a broadband thermal source such as the CMB, foam does not introduce
a spectral-line shift but instead redistributes power: the coherent-core
component is suppressed by $e^{-\sigma_\phi^2(\nu)}$ and the missing flux
reappears in a broad pedestal whose shape depends on the short-time
behaviour of the correlation function $C_\alpha(\tau)$.  Because
$\sigma_\phi^2\propto\nu^2$ (Eq.~\eqref{eq:sigmaphi}), the foam deficit
is concentrated at high frequencies while the pedestal fills in at low
frequencies, producing a spectral signature distinct from standard
chemical-potential ($\mu$) and Compton-$y$ distortions.

The Far-InfraRed Absolute Spectrophotometer (FIRAS) aboard the Cosmic
Background Explorer (COBE) measured the absolute spectrum of the CMB
monopole over $60$--$600$\,GHz by direct comparison against an internal
cryogenic blackbody calibrator --- the absolute-photometry feature that
distinguishes it from relative-photometry experiments and makes it
well suited to constraining spectral \emph{shape} distortions.  The
calibrated monopole agrees with a pure blackbody to within
$\sim 50$\,ppm (a few$\,\times10^{-5}$ fractional) across the band,
placing the still-standard limits $|\mu|\lesssim 9\times10^{-5}$ and
$|y|\lesssim 1.5\times10^{-5}$ \cite{Fixsen1996}.  These limits and the
$43$-channel calibrated spectrum form the empirical basis for the foam
constraint derived in this subsection, which sets up the broadband
signal, derives it from first principles, analyses the FIRAS
sensitivity, and contrasts the result with the GRB spectroscopic bound.

\subsubsection{Foam phase variance at CMB frequencies}
\label{sec:CMB_phasevar}

CMB photons travel from the last-scattering surface at comoving distance
$D_{\rm C}= 13.87\,{\rm Gpc}$ (Planck\,2018 flat-$\Lambda$CDM cosmology;
$H_0 = 67.36$\,km\,s$^{-1}$\,Mpc$^{-1}$, $\Omega_m=0.315$).  For foam
path-length variance accumulating uniformly in comoving distance
($d\sigma_\ell^2/dD_C$ constant, with $(1+z')^2$ retained inside the
integrand of Eq.~\eqref{eq:sigmaphi_cosmo}; anchoring~(a) of
Section~\ref{sec:differential_jitter}), the high-$z$ portion of the
path dominates the integral and the accumulated phase variance at
last scattering is amplified, tightening the FIRAS threshold to
$\alpha\lesssim 0.55$ (Section~\ref{sec:CMB_degeneracy};
$\Delta\alpha\approx +0.033$ relative to the simpler form derived
below).  Because the foam-evolution model in an expanding universe is
not specified by Eq.~\eqref{eq:NgPerlmanVar}, we treat the
cosmological enhancement as a multiplicative modification of
$A_\alpha$ that the convention $A_\alpha=1$ absorbs.

A more conservative alternative is anchoring~(b) of
Section~\ref{sec:differential_jitter}: foam variance held fixed at the
present epoch, $(1+z')^2\to 1$ inside the cosmological integrand, so
that Eq.~\eqref{eq:sigmaphi_cosmo} reduces to
\begin{equation}
  \sigma_\phi^2(\nu)
    = \left(\frac{2\pi\nu}{c}\right)^{\!2}\!
      A_\alpha\,D_{\rm C}^{2(1-\alpha)}\,\ell_P^{2\alpha},
  \label{eq:sigma_phi_CMB}
\end{equation}
where $\nu$ is the cyclic observer-frame frequency ($\omega=2\pi\nu$).
At $\nu=150$\,GHz this gives
\begin{equation}
  \sigma_\phi^2(150\,{\rm GHz})\approx
  \begin{cases}
    6.8\times10^{-2} & \alpha=1/2\ \text{(random walk)},\\
    10^{-22}         & \alpha=2/3\ \text{(holographic)}.
  \end{cases}
  \label{eq:sigma_phi_CMB_values}
\end{equation}
The strong $\nu^2$ dependence concentrates foam effects at high
frequency; for $\alpha=1/2$ the variance reaches unity at
$\nu\approx 574$\,GHz, near the top of the FIRAS band, marking where
the linearised template breaks down (Section~\ref{sec:CMB_caveats}).
The corresponding FIRAS threshold is $\alpha\lesssim 0.52$
(Section~\ref{sec:CMB_degeneracy}), which we adopt as the conservative
headline value.

\subsubsection{Foam-modified CMB spectrum}
\label{sec:CMB_derivation}

The broadband blackbody admits a direct evaluation: compute
$g^{(1)}_{\rm src}(\tau)$ as the Fourier transform of $B_\nu(T_0)$
via Eq.~\eqref{eq:WK_obs}, apply the jitter-convolution
Eq.~\eqref{eq:g1obs}, and inverse-transform to recover
$S_{\rm obs}(\nu)$.  This route entangles the model-independent
per-mode suppression $e^{-\sigma_\phi^2(\nu)}$ with the
$C_\alpha$-dependent redistribution of the suppressed flux.  We
instead present a derivation that isolates these two contributions:
the coherent core, which depends only on the asymptotic phase
variance $\sigma_\phi^2(\nu)$ and is therefore obtained in closed
form independent of foam temporal dynamics, and the pedestal kernel
$\hat{R}(\nu-\nu';\nu')$, which carries the full dependence on the
short-$\tau$ behaviour of $C_\alpha(\tau)$ and is treated
phenomenologically pending a first-principles specification of the
foam temporal autocorrelation.  A worked example with exponential
$C_\alpha(\tau)$ below makes the pedestal kernel concrete and
identifies the physically-realized case used in the Fisher analysis
of Section~\ref{sec:CMB_degeneracy}.

Decomposing the blackbody as an incoherent superposition of narrowband
modes $S_{\rm src}(\nu')=B_{\nu'}(T_0)$, each mode is modulated by the
\emph{same} achromatic jitter $\delta t(t)$, with per-mode foam suppression
\begin{equation}
  \mathcal{F}_{\rm foam}(\tau;\nu')
    = e^{-\sigma_\phi^2(\nu')[1-C_\alpha(\tau)]},
  \quad
  \sigma_\phi^2(\nu')=(2\pi\nu'/c)^2\sigma_\ell^2(D),
\end{equation}
that approaches the constant $e^{-\sigma_\phi^2(\nu')}$ for all
$\tau\neq 0$ in the $C_\alpha\to 0$ limit.

For a single narrowband mode at frequency $\nu'$, the source
coherence is $g^{(1)}_{\rm src}(\tau)=e^{-i2\pi\nu'\tau}$ (the
$\tau_G\to\infty$ limit of Eq.~\eqref{eq:g1src_gauss}), and the
jitter convolution gives
\[
  g^{(1)}_{\rm obs}(\tau)
    = e^{-i2\pi\nu'\tau}\,\mathcal{F}_{\rm foam}(\tau;\nu')
    = e^{-i2\pi\nu'\tau}\,e^{-\sigma_\phi^2(\nu')[1-C_\alpha(\tau)]}.
\]
In the dynamical limit, $\mathcal{F}_{\rm foam}(\tau;\nu')$
transitions from $\mathcal{F}_{\rm foam}(0;\nu')=1$ to the plateau
$e^{-\sigma_\phi^2(\nu')}$ within the narrow window
$|\tau|\lesssim\tau_{\rm foam}$.  The Fourier transform of the
plateau (constant in $\tau$ outside that window) is
$e^{-\sigma_\phi^2(\nu')}\delta(\nu-\nu')$ --- the coherent core at
this mode's frequency --- while the Fourier transform of the rise
from plateau to unity is supported on $|\tau|\lesssim\tau_{\rm foam}$
and is therefore broad in frequency, of width
$\sim 1/\tau_{\rm foam}$ --- the pedestal at this mode.  The split
applied below to $\mathcal{F}_{\rm foam}(\tau;\nu')$ generalises this
single-mode decomposition to the full Planck spectrum by linear
superposition.

To isolate the coherent-core
and pedestal contributions, we split
$\mathcal{F}_{\rm foam}(\tau;\nu')=e^{-\sigma_\phi^2(\nu')}+
[\mathcal{F}_{\rm foam}(\tau;\nu')-e^{-\sigma_\phi^2(\nu')}]$: the
first term contributes the constant asymptote at all lags
(Fourier-transforming to a delta function at $\nu=\nu'$), while the
residual is localised at $|\tau|\lesssim\tau_{\rm foam}$
(Fourier-transforming to a broad pedestal).  Taking the
intensity-weighted Fourier transform via Eq.~\eqref{eq:WK_obs} yields
\begin{equation}
  S_{\rm obs}(\nu)
    = e^{-\sigma_\phi^2(\nu)}\,B_\nu(T_0)
    + \int_0^\infty B_{\nu'}(T_0)\,\hat{R}(\nu-\nu';\nu')\,d\nu',
  \label{eq:CMB_spectrum}
\end{equation}
where $\hat{R}$ denotes the residual's Fourier transform.  The
coherent-core term $e^{-\sigma_\phi^2(\nu)}B_\nu(T_0)$ is exact and
model-independent (requiring only achromatic jitter and
$C_\alpha\to 0$); the pedestal kernel $\hat{R}$ depends on the
short-$\tau$ behaviour of $C_\alpha(\tau)$ --- the foam coherence
timescale $\tau_{\rm foam}$ --- and is therefore undetermined without
specifying $C_\alpha$.  This is the broadband generalisation of the
core-plus-pedestal decomposition stated in Section~\ref{sec:spec}.
Because the per-mode phase modulation $e^{i\omega\delta t(t)}$ is
unitary, any physical pedestal kernel $\hat{R}$ derived from a
specific $C_\alpha(\tau)$ conserves photon number per mode:
$\int [S_{\rm obs}(\nu)-B_\nu(T_0)]\,d\nu/(h\nu)=0$.

To make the pedestal kernel concrete, we evaluate $\hat{R}$ for the
minimal one-parameter model $C_\alpha(\tau)=e^{-|\tau|/\tau_{\rm foam}}$,
the autocorrelation of a Markovian metric fluctuation.  The framework
does not require this specific choice; any $C_\alpha(\tau)$ with
characteristic decorrelation time $\tau_{\rm foam}\ll 1/\nu_{\rm max}$
(satisfied with enormous margin for Planck-scale foam,
$\tau_{\rm foam}\sim\ell_P/c\approx 5\times10^{-44}$\,s)
has a Fourier transform of width $\gtrsim 1/\tau_{\rm foam}\gg
\nu_{\rm max}$, so the pedestal kernel is constant across the FIRAS
band irrespective of its functional form, and the conclusion that
the in-band pedestal is negligible follows from the bandwidth ratio
alone.  Substituting into Eq.~\eqref{eq:diffVarExt} gives
\begin{equation}
  \sigma_\Delta^2(\tau) = \frac{2\sigma_\ell^2(D)}{c^2}
    \bigl[1 - e^{-|\tau|/\tau_{\rm foam}}\bigr].
  \label{eq:sigmaDelta_exponential}
\end{equation}
The Fisher analysis below operates in the linearised regime
$\sigma_\phi^2(\nu')\ll 1$.  Expanding $\mathcal{F}_{\rm foam}(\tau;\nu')$
to first order and subtracting the long-lag asymptote
$e^{-\sigma_\phi^2(\nu')}$ leaves the short-$\tau$ residual
\begin{equation}
  \mathcal{F}_{\rm foam}(\tau;\nu') - e^{-\sigma_\phi^2(\nu')}
  \;\approx\; \sigma_\phi^2(\nu')\,e^{-|\tau|/\tau_{\rm foam}},
  \label{eq:Ffoam_residual_exponential}
\end{equation}
whose Fourier transform is a Lorentzian centred on the source
frequency,
\begin{equation}
  \hat{R}(\nu-\nu';\nu')
    \;=\; \sigma_\phi^2(\nu')\,
    \frac{2\tau_{\rm foam}}{1 + [2\pi(\nu-\nu')\tau_{\rm foam}]^2},
  \label{eq:R_lorentzian}
\end{equation}
of full width at half maximum $1/(\pi\tau_{\rm foam})$.  For Planck-scale foam
$\tau_{\rm foam}\sim\ell_P/c\approx 5\times 10^{-44}$\,s the
Lorentzian width is $\sim 6\times 10^{42}$\,Hz, eleven orders of
magnitude beyond the FIRAS upper edge $\nu_{\rm max}=6\times 10^{11}$\,Hz;
the fraction of the redistributed flux falling within the FIRAS band
is $\sim 10^{-31}$.  The physical pedestal therefore deposits no
detectable in-band power for any plausible $\tau_{\rm foam}$.  The
Fisher analysis of Section~\ref{sec:CMB_degeneracy} consequently
adopts the no-pedestal limit as the primary template.

To probe robustness against in-band pedestal contamination beyond
the physical-kernel result above, we also evaluate the Fisher matrix
with an \emph{adversarial} band-filling redistribution kernel,
deliberately shaped $\propto 1/\nu$ to match the Rayleigh--Jeans form
of the $\mu$-distortion template (Section~\ref{sec:CMB_degeneracy}) and
thereby maximize the foam--$\mu$ Fisher degeneracy --- the most
effective way an in-band pedestal could disguise foam as a standard
$\mu$-distortion:
\begin{equation}
  K_\nu^{\rm fast}(\nu) \;=\;
  \begin{cases}
    \dfrac{1}{\nu\,\ln(\nu_{\rm max}/\nu_{\rm min})} &
    \nu_{\rm min}\le\nu\le\nu_{\rm max},\\[4pt]
    0 & \text{otherwise},
  \end{cases}
  \label{eq:Knufast}
\end{equation}
with $\nu_{\rm min}=60$\,GHz and $\nu_{\rm max}=600$\,GHz set by the
FIRAS band edges.  $K_\nu^{\rm fast}(\nu)$ is a normalized density
in observer frequency ($\int_{\nu_{\rm min}}^{\nu_{\rm max}}
K_\nu^{\rm fast}\,d\nu=1$, units Hz$^{-1}$) that specifies the
fraction of the total foam-removed flux reappearing as pedestal
intensity per Hz at observer frequency $\nu$; the form above
corresponds to uniform energy redistribution in $\ln\nu$ across the
band.  Within the FIRAS band, $K_\nu^{\rm fast}(\nu)$ replaces the
physical pedestal contribution to Eq.~\eqref{eq:CMB_spectrum} by an
adversarial factorised substitution
\begin{equation}
  \int B_{\nu'}(T_0)\,\hat{R}(\nu-\nu';\nu')\,d\nu'
    \;\longrightarrow\;
    \sigma_\ell^2(D)\,K_\nu^{\rm fast}(\nu)\,P_{\rm foam},
  \quad
  P_{\rm foam} \equiv \int_{\nu_{\rm min}}^{\nu_{\rm max}}
    \!\left(\frac{2\pi\nu'}{c}\right)^{\!2} B_{\nu'}(T_0)\,d\nu',
  \label{eq:Kfast_pedestal}
\end{equation}
with $P_{\rm foam}$ the total foam-removed flux per unit
$\sigma_\ell^2$.  The substitution forces the pedestal kernel to
factorise into a $\nu$-only template ($K_\nu^{\rm fast}$) times a
$\nu'$-only weighting, a structure no physical
$\hat{R}(\nu-\nu';\nu')$ satisfies (the Lorentzian above carries
$\nu'$-dependence through $\sigma_\phi^2(\nu')$).  The template is
energy-conserving but not photon-number-conserving, and serves purely
as an adversarial probe of the foam--$\mu$ kernel-shape degeneracy.
As Section~\ref{sec:CMB_degeneracy} shows, this adversarial choice does
not in fact loosen the bound: by depositing an in-band pedestal it
distorts the spectrum away from a pure blackbody and so makes the foam
\emph{more} detectable, leaving the no-pedestal deficit as the
conservative case.

\subsubsection{Fisher analysis}
\label{sec:CMB_degeneracy}

We use $I_\nu\equiv S_{\rm obs}(\nu)$ to denote the observed CMB
intensity including the foam signal.
Working in the linearised regime
$\sigma_\phi^2(\nu)\ll 1$, we treat four spectral parameters
$\boldsymbol{\theta}=(A_{\rm eff},T_0,\mu,y)$, where
$A_{\rm eff}=A_\alpha D_C^{2(1-\alpha)}\ell_P^{2\alpha}=\sigma_\ell^2(D_C)$
is the cumulative foam path-length variance at the CMB distance
(Eq.~\eqref{eq:NgPerlmanVar}), so that
$\sigma_\phi^2(\nu)=(2\pi\nu/c)^2 A_{\rm eff}$.  Note that $\alpha$ is
not fitted directly; only the lumped amplitude $A_{\rm eff}$ enters
the Fisher matrix, and bounds on $\alpha$ are obtained
\emph{a posteriori} by inverting the Ng--Perlman relation at fixed
$A_\alpha=1$.

In the physically-realized regime (Section~\ref{sec:CMB_derivation})
the redistributed pedestal lies at $\sim 10^{42}$\,Hz, eleven orders of
magnitude above the FIRAS band, so the only in-band foam signature is
the pure deficit
\begin{equation}
  \frac{\partial I_\nu}{\partial A_{\rm eff}}
    = -\!\left(\frac{2\pi\nu}{c}\right)^{\!2} B_\nu(T_0),
  \label{eq:dIdA_phys}
\end{equation}
a $\nu^2$-weighted (blue-tilted) Planck spectrum.  The
chemical-potential and Compton-$y$ distortion templates are
\begin{align}
  \frac{\partial I_\nu}{\partial\mu}
    &= -B_\nu(T_0)\,\frac{x e^x}{(e^x-1)^2}, \label{eq:dIdmu}\\
  \frac{\partial I_\nu}{\partial y}
    &= \phantom{-}B_\nu(T_0)\,\frac{x e^x}{e^x-1}
       \!\left[\frac{x(e^x+1)}{e^x-1}-4\right], \label{eq:dIdy}
\end{align}
with $x=h\nu/(k_B T_0)$, encoding respectively early-universe energy
injection (Silk damping, decaying particles) and late-time electron
heating (reionization, cluster gas).

A Fisher analysis using the 43 FIRAS channels~\cite{Fixsen1996} with
diagonal noise approximated as a flat $\sigma_\nu=0.01$\,MJy/sr per
channel (a band-averaged simplification of the published
channel-by-channel uncertainties) quantifies how well $A_{\rm eff}$ can
be separated from these templates.  Raising $T_0$ brightens the whole
Planck curve, while the foam deficit suppresses it with an extra
$\nu^2$ weighting: the two differ only by a spectral tilt that the
CMB-bright band (the peak through the onset of the Wien fall-off)
barely resolves.  The foam template is therefore degenerate with a
temperature shift at $\rho_{A_{\rm eff}T_0}=0.9997$, and through that
near-degeneracy it is also strongly collinear with the $y$ and $\mu$
templates ($\rho_{A_{\rm eff}y}=0.999$, $\rho_{A_{\rm eff}\mu}=0.93$).
The marginal amplitude uncertainty is correspondingly large,
\begin{equation}
  \sigma(A_{\rm eff}) = 2.2\times10^{-11}\,{\rm m}^2,
  \qquad
  \sigma(A_{\rm eff})\big|_{\mu,y\,\text{fixed}}
    = 8.2\times10^{-13}\,{\rm m}^2,
  \label{eq:sigma_Aeff}
\end{equation}
and setting $A_{\rm eff}^{\rm Ng\text{-}Perlman}(\alpha,D_C)
=\sigma(A_{\rm eff})$ (Fig.~\ref{fig:cmb_snr}) gives the conservative
headline threshold $\alpha\lesssim 0.52$.

This is a best case.  The noise model is foreground-free, whereas
galactic thermal dust and the cosmic infrared background dominate the
Wien side of the spectrum --- precisely where the $\nu^2$-tilt would
finally separate the foam deficit from a pure temperature shift.
Marginalising over those foreground components would tighten the
foam--$T_0$ degeneracy further and weaken the bound, so
$\alpha\lesssim 0.52$ is the most optimistic FIRAS reach rather than a
realistic forecast.

\begin{figure}[h]
\centering
\includegraphics[width=0.82\columnwidth]{cmb_snr_alpha}
\caption{Detection SNR for the foam amplitude $A_{\rm eff}$ as a function
of foam exponent $\alpha$ ($A_\alpha=1$, $D_C=13.87$\,Gpc).
Blue curves: marginal SNR (all four parameters $A_{\rm eff},T_0,\mu,y$ free).
Red curves: conditional SNR ($\mu,y$ held fixed), which is $\sim 27\times$ higher.
Line styles denote the instrument: solid FIRAS, dashed PIXIE,
dash-dotted SuperPIXIE, dotted Voyage~2050, with forecast
$\mu$-sensitivities scaled relative to FIRAS \cite{Chluba2021}.
The marginal thresholds (SNR\,$=1$) cluster at $\alpha\approx 0.52$--$0.56$,
illustrating the logarithmic saturation of the bound with sensitivity.
The dotted horizontal line marks SNR\,$=1$; vertical dashed lines mark the
spectroscopic lower bound from GRB\,221009A ($\alpha\gtrsim 0.63$) and the
holographic prediction $\alpha=2/3$.}
\label{fig:cmb_snr}
\end{figure}

To bound the impact of adversarial in-band pedestal contamination, we
repeat the analysis with the adversarial band-filling kernel of
Section~\ref{sec:CMB_derivation}, replacing the deficit by
\begin{equation}
  \frac{\partial I_\nu}{\partial A_{\rm eff}}
    = -\!\left(\frac{2\pi\nu}{c}\right)^{\!2} B_\nu(T_0)
      + K_\nu^{\rm fast}(\nu)\!\int_{\nu_{\rm min}}^{\nu_{\rm max}}
        \left(\frac{2\pi\nu'}{c}\right)^{\!2} B_{\nu'}(T_0)\,d\nu',
  \label{eq:dIdA}
\end{equation}
whose $1/\nu$ pedestal we deliberately gave the Rayleigh--Jeans shape
of the $\mu$-distortion template (Eq.~\eqref{eq:Knufast}; the
Rayleigh--Jeans limit of Eq.~\eqref{eq:dIdmu} is
$\Delta I/B\to-\mu/x\propto\nu^{-1}$) so as to maximise the foam--$\mu$
overlap.  The resulting correlation matrix for
$(A_{\rm eff},\,T_0,\,\mu,\,y)$ is
\begin{equation}
\rho_{ij} = \begin{pmatrix}
  1.000 & 0.937 & \mathbf{0.976} & 0.010 \\
  0.937 & 1.000 & 0.969 & -0.212 \\
  \mathbf{0.976} & 0.969 & 1.000 & -0.127 \\
  0.010 & -0.212 & -0.127 & 1.000
\end{pmatrix}.
\label{eq:fisher_corr_matrix}
\end{equation}
Adding the in-band pedestal rotates the foam template off the
temperature axis ($\rho_{A_{\rm eff}T_0}:0.9997\to0.937$) and off the
$y$ template ($\rho_{A_{\rm eff}y}:0.999\to0.010$), leaving the
engineered foam--$\mu$ degeneracy ($\rho_{A_{\rm eff}\mu}=0.976$, bold)
as the dominant one.  Because $0.976$ is milder than the $0.9997$ it
replaces, the amplitude is in fact \emph{better} constrained,
\begin{equation}
  \sigma(A_{\rm eff}) = 1.7\times10^{-12}\,{\rm m}^2,
  \qquad
  \sigma(A_{\rm eff})\big|_{\mu,y\,\text{fixed}}
    = 2.9\times10^{-13}\,{\rm m}^2,
  \label{eq:sigma_Aeff_adv}
\end{equation}
a factor $\sim 6$ marginalisation penalty driven by the foam--$\mu$
collinearity, and the threshold \emph{tightens} slightly to
$\alpha\lesssim 0.53$.  The adversarial in-band pedestal therefore does
not loosen the conservative bound; it shifts it by only
$\Delta\alpha\approx 0.009$, the near-coincidence reflecting the
logarithmic Ng--Perlman relation
$\sigma_\ell^2 = D_C^{2(1-\alpha)}\ell_P^{2\alpha}$: at fixed
$\sigma_\ell$, $\Delta\alpha=\ln N/[2\ln(D_C/\ell_P)]\approx 0.009$ for
the factor $N\approx 13$ separating the two amplitude uncertainties.
An external constraint on $\mu$ helps most in this $\mu$-limited worst
case: imposing the $\Lambda$CDM dissipation prediction
$\mu\approx 2\times10^{-8}$~\cite{Chluba2021} as a prior (five orders
of magnitude tighter than the FIRAS $\mu$ measurement, hence
effectively fixing $\mu$) sharpens $\sigma(A_{\rm eff})$ by
$\sim 3$--$5\times$, though the logarithmic relation again limits the
threshold gain to $\Delta\alpha\lesssim 0.005$; in the physical
no-pedestal case, where the limiting degeneracy is with $T_0$ rather
than $\mu$, even fixing $\mu$ barely helps.  Either way, the worked
example of Section~\ref{sec:CMB_derivation} shows that any
exponential-class $C_\alpha$ with Planck-scale $\tau_{\rm foam}$ lies
in the no-pedestal regime, so the physical $\alpha\lesssim 0.52$ is the
headline and the adversarial $\alpha\lesssim 0.53$ merely confirms its
robustness.

The proposed PIXIE satellite~\cite{Kogut2011} --- a fully cryogenic
($T_{\rm sys}\approx2.7$\,K) nulling polarizing Fourier transform spectrometer
spanning $30$\,GHz to $6$\,THz --- is projected to surpass FIRAS sensitivity
by roughly three orders of magnitude, reaching $|\mu|\lesssim10^{-8}$ and
$|y|\lesssim10^{-9}$ and entering the regime of guaranteed $\Lambda$CDM
spectral signals from Silk damping and reionization.
Propagating this gain through the Fisher-matrix scaling derived above,
the marginal foam detection threshold shifts from
$\alpha\lesssim0.52$ (FIRAS) to approximately $\alpha\lesssim0.55$ (PIXIE),
narrowing the gap with the spectroscopic lower
bound $\alpha\gtrsim0.63$ from GRB\,221009A (Fig.~\ref{fig:cmb_snr}).
Both thresholds are quoted for the physical no-pedestal deficit and
the observer-frame $\sigma_\phi^2$ of Eq.~\eqref{eq:sigmaphi}; the
cosmological-redshift weighting of Eq.~\eqref{eq:sigmaphi_cosmo} could
shift each by up to $\Delta\alpha\sim 0.03$.
The limiting degeneracy --- with a temperature shift for the
no-pedestal deficit --- persists on the linearized template, but
PIXIE's extension to $6$\,THz accesses the nonlinear regime
$\sigma_\phi^2\gg1$ for $\alpha\lesssim0.55$ under the
cosmological-redshift weighting of Eq.~\eqref{eq:sigmaphi_cosmo}
(observer-frame Eq.~\eqref{eq:sigmaphi} shifts this to
$\alpha\lesssim 0.52$), where the foam deficit diverges markedly from
both the blackbody and $\mu$-distortion shapes and the degeneracy may
be substantially reduced.  The $\alpha\lesssim0.55$ projection is
therefore conservative: if the nonlinear template effectively breaks
this degeneracy at PIXIE sensitivity, the true threshold could be
significantly more stringent.

\subsubsection{Bias in recovered $T_0$ from unmodeled foam}
\label{sec:CMB_T0_bias}

The Fisher matrix also quantifies how much the fitted CMB temperature
would shift if foam were present but omitted from the model.  We
evaluate this for the physically-realized no-pedestal deficit
(Eq.~\eqref{eq:dIdA_phys}).  Suppose the data contain foam at amplitude
$A_{\rm eff}^{\rm true}$ but the analyst fits only
$\boldsymbol{\theta}_{\rm fit}=(T_0,\mu,y)$.  In the linearised Fisher
regime the parameter shift is
\begin{equation}
  \hat{\boldsymbol{\theta}}_{\rm fit}
    - \boldsymbol{\theta}_{\rm fit}^{\rm true}
  = -(F_{\rm fit,fit})^{-1}\,F_{\rm fit,omit}\,A_{\rm eff}^{\rm true},
  \label{eq:bias_formula}
\end{equation}
where $F_{\rm fit,fit}$ is the $3\times3$ Fisher block for
$(T_0,\mu,y)$ and $F_{\rm fit,omit}$ couples the fitted parameters to
the omitted foam amplitude, both evaluated for the no-pedestal template
(Eqs.~\eqref{eq:dIdA_phys}--\eqref{eq:dIdy}).

The $T_0$ component gives
\begin{equation}
  \Delta T_0 = b_{T_0}\,A_{\rm eff}^{\rm true},
  \qquad b_{T_0} = -\bigl[(F_{\rm fit,fit})^{-1}F_{\rm fit,omit}\bigr]_{T_0}.
  \label{eq:T0_bias}
\end{equation}
Numerically, $b_{T_0}=+1.5\times10^7\,{\rm K\,m^{-2}}$: because the
$\nu^2$-tilted deficit is nearly a temperature shift
($\rho_{A_{\rm eff}T_0}=0.9997$), unmodeled foam is absorbed largely
into $T_0$, biasing the recovered temperature high.  (Unlike the
adversarial pedestal, whose Rayleigh--Jeans excess the $\mu$ template
soaks up, $b_\mu\approx 1.0\times10^8\,{\rm m^{-2}}$, the no-pedestal
deficit carries no such $\mu$-like component, $b_\mu\approx
4.5\times10^6\,{\rm m^{-2}}$, and so projects directly onto $T_0$.)
Requiring $|\Delta T_0|<\sigma(T_0)=5.7\times 10^{-4}$\,K
\cite{Fixsen2009} gives
\begin{equation}
  A_{\rm eff}^{\rm max,\,bias}
    = \frac{\sigma(T_0)}{|b_{T_0}|} = 3.9\times 10^{-11}\,{\rm m}^2,
  \label{eq:Aeff_bias_bound}
\end{equation}
corresponding (with $A_\alpha=1$) to a crossover exponent
$\alpha_{\rm bias}=0.518$.  Models with $\alpha<0.518$ would have
biased the Fixsen $T_0$ by more than $1\sigma$ had such a signal been
present and unmodeled; models with $\alpha>0.518$ would not.

The crossover $\alpha_{\rm bias}\approx 0.518$ essentially coincides
with the marginal detection threshold $\alpha\lesssim 0.52$, and the
coincidence is physical: the same foam--$T_0$ near-degeneracy that
limits detection (Section~\ref{sec:CMB_degeneracy}) also channels
unmodeled foam into the temperature.  The deficit is hard to
distinguish from a temperature shift, so it is simultaneously difficult
to \emph{detect} and, if omitted, efficiently absorbed into $T_0$;
detection threshold and $T_0$-bias crossover are therefore set by the
same degeneracy.  Both probe the regime below the GRB lower exclusion
bound $\alpha\gtrsim 0.63$ --- FIRAS sets an upper detectability limit
while the GRB analysis sets a lower exclusion bound.
Figure~\ref{fig:cmb_bias} plots $|\Delta T_0|/\sigma(T_0)$ as a
function of $\alpha$; at the detection threshold $\alpha\approx 0.52$
(marginal SNR $\approx 1$) the ratio is $\approx 0.6$, falling below
$10^{-2}$ only for $\alpha\gtrsim 0.54$.  The published Fixsen
$T_0=2.72548$\,K is therefore unbiased at the $\ll 10^{-4}\sigma$ level
by any foam model consistent with the GRB bound $\alpha\gtrsim 0.63$,
although a hypothetical strong-foam model at the very edge of FIRAS
detectability ($\alpha\approx 0.52$) would shift it by up to
$\sim 0.6\,\sigma(T_0)$.

\begin{figure}[h]
\centering
\includegraphics[width=0.82\columnwidth]{cmb_bias_alpha}
\caption{Foam-induced bias in the recovered CMB temperature,
$|\Delta T_0|/\sigma(T_0)$, as a function of foam exponent $\alpha$
when foam at amplitude $A_{\rm eff}$ ($A_\alpha=1$, $D_C=13.87$\,Gpc)
is present but omitted from the spectral fit.
Solid: FIRAS with $\sigma(T_0)=5.7\times10^{-4}$\,K~\cite{Fixsen2009};
dashed: PIXIE with $\sigma(T_0)\approx 5.7\times10^{-7}$\,K (noise floor
$\sim 10^3\times$ lower).
The dotted horizontal line marks $|\Delta T_0|=\sigma(T_0)$; the FIRAS
crossover occurs at $\alpha\approx 0.518$.
Vertical dashed lines mark the spectroscopic lower bounds and the
holographic prediction $\alpha=2/3$ (as in Fig.~\ref{fig:cmb_snr}).}
\label{fig:cmb_bias}
\end{figure}

\subsubsection{Caveats: nonlinear regime and parameter degeneracy}
\label{sec:CMB_caveats}

Two caveats qualify the linearised analysis.

\emph{Nonlinear regime.}  The linearisation
$\Delta I_\nu\approx\sigma_\phi^2(\nu)\times[\text{template}]$ is valid
only for $\sigma_\phi^2\ll 1$.  For $\alpha=1/2$, $\sigma_\phi^2=1$ is
reached at $\nu\approx 574$\,GHz, within the FIRAS band.  Replacing
the linearised deficit $-\sigma_\phi^2(\nu)B_\nu$ by the saturated form
$-(1-e^{-\sigma_\phi^2(\nu)})B_\nu$ reduces the \emph{integrated}
in-band foam amplitude (summed over FIRAS channels) by $\approx 13\%$
at $\alpha=0.50$ and by $<1\%$ at $\alpha\gtrsim 0.52$.
At the headline threshold $\alpha\lesssim 0.52$ the saturation
correction is thus only $\approx 1\%$
($\Delta\alpha\sim 3\times 10^{-5}$), so the linearised bound is
unaffected at the $\Delta\alpha\ll 10^{-3}$ level; the nonlinear
template becomes important only for stronger foam
($\alpha\lesssim 0.50$), where it may also help break the foam--$\mu$
degeneracy.

\emph{Parameter degeneracy within the foam template.}  The exponent
$\alpha$ and the amplitude $A_\alpha$ enter only through the product
$A_\alpha D_C^{2(1-\alpha)}\ell_P^{2\alpha}=A_{\rm eff}$.  FIRAS
constrains this product directly but not the exponent separately;
converting a bound on $A_{\rm eff}$ to a bound on $\alpha$ requires
fixing $A_\alpha$ by convention (here $A_\alpha=1$, following Ng \&
Perlman~\cite{NgPerlman2022}).  Any quoted bound on $\alpha$ carries
this implicit normalisation.

\subsubsection{Complementarity with GRB\,221009A}
\label{sec:CMB_complementarity}

The FIRAS sensitivity range $\alpha\lesssim 0.52$--$0.55$ does not
overlap the GRB spectroscopic lower bound $\alpha\gtrsim 0.63$ of
Section~\ref{sec:grb221009a}: FIRAS probes strictly lower exponents.
Any foam model consistent with the GRB bound has $\sigma_\phi^2\ll 1$
at CMB frequencies, so the $T_0$ bias is entirely negligible
(Section~\ref{sec:CMB_T0_bias} and Fig.~\ref{fig:cmb_bias}).

FIRAS provides an independent broadband test of foam --- absolute
calibrated photometry, no source-profile assumption, and a different
systematic-error budget from spectroscopy --- and constrains the
parameter space below the GRB bound where more exotic foam models
would reside.  A definitive FIRAS foam analysis would require
specifying $C_\alpha(\tau)$ to determine the physical pedestal shape
and replacing the linearised template with the full nonlinear
expression at high frequencies where $\sigma_\phi^2\gtrsim 1$.
The next generation of absolute spectrophotometers would deepen this
reach only marginally.  Extending the Fisher analysis above to the
$10$\,GHz--$3$\,THz band of the SuperPIXIE and Voyage~2050 concepts
\cite{Chluba2021} --- with their forecast distortion sensitivities
$\sigma(\mu)\approx 7.7\times 10^{-9}$ and $2.2\times 10^{-9}$,
respectively --- shifts the marginal foam threshold only to
$\alpha\lesssim 0.55$ (SuperPIXIE) and $\alpha\lesssim 0.56$
(Voyage~2050), barely beyond the PIXIE projection.  The reason is
structural rather than instrumental: in the physically-realized
no-pedestal limit the foam deficit $-\sigma_\phi^2(\nu)B_\nu$ is
degenerate with a temperature shift at $\rho_{A T_0}=0.9997$, and this
degeneracy is \emph{band-robust} --- broadening the coverage from the
FIRAS band to $10$\,GHz--$6$\,THz leaves it unchanged to four decimal
places and even slightly worsens the foam-to-$\mu$ sensitivity ratio
$\sigma(A_{\rm eff})/\sigma(\mu)$.  The wider band does break the
foam--$\mu$ degeneracy probed by the adversarial $K_\nu^{\rm fast}$
pedestal (from $\rho_{A\mu}=0.98$ at FIRAS to $0.66$ over
$10$\,GHz--$3$\,THz), so future missions are markedly more robust
against adversarial in-band pedestal contamination; but this does not
move the physical threshold.  Together with the logarithmic
Ng--Perlman scaling
$\Delta\alpha=\ln N/[2\ln(D_C/\ell_P)]$, which converts even a
$10^{4}$-fold sensitivity gain over FIRAS into $\Delta\alpha\lesssim
0.04$, this places a firm ceiling: spectral-distortion photometry
probes the foam parameter space at $\alpha\lesssim 0.56$ and cannot,
by sensitivity or bandwidth alone, reach the GRB spectroscopic bound
$\alpha\gtrsim 0.63$ or the holographic value $\alpha=2/3$.  The two
channels are therefore permanently complementary rather than
overlapping.

% ==============================================================
\section{Conclusions}
\label{sec:conclusions}

All single-source spectroscopic bounds derived in this paper assume the
dynamical-foam limit $\tau_{\rm foam}\ll\tau_c$
(Section~\ref{sec:closedforms}), in which the foam path-length jitter
decorrelates many times within each photon coherence time.  If foam
fluctuations are instead slow ($\tau_{\rm foam}\gg\tau_c$), the
accumulated phase is essentially frozen during each coherence time and
the foam acts as a deterministic, frequency-independent phase screen
--- no spectral broadening or peak reduction occurs within a single
observation, and the shape-distortion bounds do not apply.  The CMB
constraint (Section~\ref{sec:CMB}) is similarly derived in the
dynamical limit (Section~\ref{sec:CMB_derivation} works in
$C_\alpha\to 0$): in the opposite quasi-static limit the per-mode
foam suppression $\mathcal{F}_{\rm foam}(\tau)\to 1$
(Eq.~\eqref{eq:Ffoam}) and the CMB monopole is unmodified, so the
shape-distortion bounds do not apply in either channel.

The dynamical bound $\alpha\gtrsim 0.63$
(Eq.~\eqref{eq:alpha_grb_dyn}) does not exclude the holographic value
$\alpha=2/3\approx 0.667$.  A principal instrumental limitation of the
present analysis is the MeV sensitivity gap: the spectral region
between the upper edge of the Fermi-GBM BGO detector response
($\lesssim40\,\mathrm{MeV}$, where the effective area falls steeply)
and the regime in which Fermi-LAT achieves robust detection efficiency
($\gtrsim100\,\mathrm{MeV}$).
The frequency $\nu_\star \simeq 37\,\mathrm{MeV}$ at which the foam-induced
spectral break is predicted to be most constraining falls squarely within
this gap, so neither instrument provides fully resolved coverage across
the critical transition.
Closing the MeV gap is a recognized community priority, addressed by
next-generation satellite concepts that combine Compton-scattering and
pair-conversion detection techniques to achieve continuous sensitivity from
$\sim100\,\mathrm{keV}$ to beyond $1\,\mathrm{GeV}$.
The All-sky Medium Energy Gamma-ray Observatory
(AMEGO~\cite{McEnery2019}) and its MIDEX-scale variant
AMEGO-X~\cite{Caputo2022} are designed to cover $100\,\mathrm{keV}$ to
$1\,\mathrm{GeV}$ with order-of-magnitude improvements over both COMPTEL
and Fermi-LAT in the MeV decade; the ESA-proposed M-class concept
e-ASTROGAM~\cite{DeAngelis2017} targets a comparable bandpass.
The near-term NASA SMEX mission COSI (scheduled launch $\sim$2027)
will extend broadband sensitivity to $0.2$--$5\,\mathrm{MeV}$, providing
a valuable pathfinder, though its upper energy bound falls well short of
the $37\,\mathrm{MeV}$ regime and the gap therefore remains open.
Realization of any of the broader-bandpass concepts would, for the first
time, permit fully resolved spectral measurements of GRB prompt emission
across the $37\,\mathrm{MeV}$ transition without the instrumental
discontinuity that limits the present work; reducing the per-epoch
line-width uncertainties to $\delta r\lesssim 0.01$ would tighten the
bound to $\alpha\gtrsim 0.64$, removing a major systematic and
enabling substantially tighter tests of spacetime-foam models.

The multi-epoch $r_{\rm obs}$ dataset constrains any propagation
effect that broadens an emission line in proportion to a power of
$E$, not only the Ng--Perlman spacetime-foam template considered
here.  The most natural extension is to stochastic Lorentz-invariance
violation~\cite{AmelinoCamelia2013}, in which the photon group
velocity carries a Planck-suppressed random coefficient
$v(E)=c\bigl[1-\xi\,(E/E_{\rm QG,n})^n\bigr]$, with $\xi$ of order
unity drawn independently for each Planck-scale traversal.  Photons
in the same line then propagate at slightly different speeds,
broadening the observed width as
$\sigma_{\rm LIV}(E)/E\propto E^n$ on top of any intrinsic
contribution.  Unlike the deterministic LIV signal probed by photon
time-of-flight studies --- where GRB\,221009A LHAASO TeV photons
already constrain $E_{\rm QG,1}\gtrsim$ a
few~$\times\,E_{\rm Pl}$~\cite{LHAASO2024} --- stochastic LIV
produces a \emph{width} signal that survives light-curve
cross-correlation, and a multi-epoch $r_{\rm obs}$ measurement is
therefore one of the few clean probes available.  Fitting
Eq.~\eqref{eq:robs_model} recovers the foam constraint as the
$n=1$ case; refitting the same seven-epoch dataset under
$r_{\rm LIV}\propto(E/E_0)^n$ would yield bounds on
$\langle\xi^2\rangle^{1/2}/E_{\rm QG,n}^n$ at each order, with the
$n=2$ (quadratic) case particularly valuable because the
deterministic time-of-flight tests are weakest there at $z=0.151$.

\begin{acknowledgments}
This work was supported in part by the Director, Office of Science, Office of
High Energy Physics of the U.S.\ Department of Energy under Contract No.\
DE-AC02-05CH11231.

The authors used Claude (Anthropic) and Valency Systems to assist with drafting and editing portions of this manuscript, with detailed reference searches, and in developing the Fisher-matrix and GRB multi-epoch analysis scripts. All content, equations, numerical results, and conclusions were reviewed and verified by the authors, who take full responsibility for the work.  AI assistance was used in accordance with \textit{Physical Review} editorial policy on computational tools.

We thank Robin Kaiser for helpful discussions that helped motivate this work.
\end{acknowledgments}

\bibliographystyle{apsrev4-2}
\bibliography{../references}

\end{document}
